% Faithful transcription — original Italian (18th c.). Transcribed from the 1911 Opere
% Matematiche, vol. II (Wikimedia Commons djvu), the complete memoir "Metodo per misurare la
% lemniscata, Schediasma II", printed pages 304–313 (= djvu pages 334–343). \origpage uses the
% PRINTED page numbers. Content and archaic spelling follow the print; only markup is
% standardized (PLAN.md §4.4). NOTE: the heading on p. 308 prints "Teorema IV" a second time for
% what every later cross-reference in this same memoir calls "il teorema VI" (pp. 309, 310) —
% almost certainly a printer's repeat of the earlier "Teorema IV" heading (p. 305) rather than
% the intended "Teorema VI". Reproduced faithfully as printed, not renumbered. Likewise p. 305's
% "Teorema IV" heading itself cites "equazioni (6), e (8)" where every other theorem in this
% pattern cites the two equations it is about to state — here that would be (7) and (8); "(6)" is
% kept as printed (probable misprint) since eq. (6) belongs to the preceding Teorema III. The
% "Terzo" remark in Scolio I (p. 306) prints $\int dz/\sqrt{1+x^4}$ (mixing differential $dz$ with
% the variable $x$ under the radical); given its consequence in the "Sesto" remark reads
% $\int dx/\sqrt{1+x^4}$ throughout, "dz" is very likely a compositor's error for "dx" — kept as
% printed. Fig. 32, printed on the same plate (Tav. III) as figs. 29, 30, 31, 33, is not cited by
% number anywhere in this memoir's text and is therefore not reproduced here.
\documentclass{article}
\usepackage{readmasters}
\begin{document}

\origpage{304}
\section*{XXXIV. Metodo per misurare la lemniscata}
\subsection*{Schediasma II (*)}
\emph{(*) Giornale de' letterati d'Italia, tomo XXX, pag.~87.}

\textbf{Teorema I.} --- Sieno le due infrascritte equazioni (1), e (2), io dico, che posta la
prima di esse, sussiste anche l'altra

\[ x = \frac{\sqrt{1 \mp \sqrt{1-z^{4}}}}{z}. \tag{1} \]

\[ \frac{\pm dz}{\sqrt{1-z^{4}}} = \frac{dx\sqrt{2}}{\sqrt{1+x^{4}}}. \tag{2} \]

\textbf{Dimostrazione.} --- L'equazione (1) differenziata somministra un valore di $dx$,
ch'essendo ridotto ad una semplice espressione fa conoscere

\[ dx = \frac{\pm dz\sqrt{1 \mp \sqrt{1-z^{4}}}}{z^{2}\sqrt{1-z^{4}}}. \]

Dalla medesima equazione (1) si deduce ancora la seguente

\[ \frac{\sqrt{1+x^{4}}}{\sqrt{2}} = \frac{\sqrt{1 \mp \sqrt{1-z^{4}}}}{z^{2}}. \]

E dividendo la penultima equazione per quest'ultima si giunge all'equazione (2).

Il che dovea dimostrarsi.

\textbf{Teorema II.} --- Sieno le due equazioni infrascritte (3), e (4); io dico, che posta la
prima di essa sussiste anche l'altra

\[ x = \frac{\sqrt{1\mp z}}{\sqrt{1\pm z}}. \tag{3} \]

\[ \frac{\mp dz}{\sqrt{1-z^{4}}} = \frac{dx\sqrt{2}}{\sqrt{1+x^{4}}}. \tag{4} \]

\origpage{305}
\textbf{Dimostrazione.} --- Differenziando l'equazione (3), e operando nel debito modo, si à
$dx = \dfrac{\pm dz}{\sqrt{1-z^{2}}} \times \dfrac{1}{1\pm z}$.

Inoltre la stessa equazione (3) mostra $\dfrac{\sqrt{1+x^{4}}}{\sqrt{2}} = \dfrac{\sqrt{1+z^{2}}}{1\pm z}$.

E dividendo la penultima equazione per quest'ultima si ottiene la equazione (4).

Il che dovea dimostrarsi.

\textbf{Teorema III.} --- Sieno le due equazioni infrascritte (5), e (6); io dico, che posta la
prima di esse, sussiste anche l'altra

\[ x = \frac{u\sqrt{2}}{\sqrt{1-u^{4}}}. \tag{5} \]

\[ \frac{du}{\sqrt{1-u^{4}}} = \frac{1}{\sqrt{2}}\times\frac{dx}{\sqrt{1+x^{4}}}. \tag{6} \]

\textbf{Dimostrazione.} --- Dall'equazione (5) differenziata, e maneggiata a dovere risulta
$\dfrac{dx}{\sqrt{2}} = \dfrac{du}{\sqrt{1-u^{4}}} \times \dfrac{1+u^{4}}{1-u^{4}}$ e dalla
suddetta equazione (5) si deduce $\sqrt{1+x^{4}} = \dfrac{1+u^{4}}{1-u^{4}}$.

Dividendo poscia la penultima equazione per quest'ultima, si arriva all'equazione (6).

Il che dovea dimostrarsi.

\textbf{Teorema IV.} --- Sieno le due infrascritte equazioni (6), e (8); io dico, che posta la
prima di esse, sussiste anche l'altra

\[ x = \frac{\sqrt{1-t^{4}}}{t\sqrt{2}}. \tag{7} \]

\[ \frac{-dt}{\sqrt{1-t^{4}}} = \frac{1}{\sqrt{2}}\times\frac{dx}{\sqrt{1+x^{4}}}. \tag{8} \]

\textbf{Dimostrazione.} --- Differenziando l'equazione (7), e riducendo il valore di $dx$ ad una
comoda espressione, ritrovasi $dx\sqrt{2} = \dfrac{-dt}{\sqrt{1-t^{4}}} \times \dfrac{1+t^{4}}{t^{2}}$.

La stessa equazione (7) fa vedere $2\sqrt{1+x^{4}} = \dfrac{1+t^{4}}{t^{2}}$.

\origpage{306}
E la penultima equazione divisa per quest'ultima rende l'equazione (8).

Il che dovea dimostrarsi.

\textbf{Scolio I.} --- Per far uso di questi teoremi si noti:

Primo, che $\displaystyle\int \frac{dz}{\sqrt{1-z^{4}}}$ rappresenta l'arco diretto $CS$ della
lemniscata (fig.~29) sostenuto dalla corda $CS = z$, purchè il semiasse $CL$ di questa curva sia
eguale all'unità.

\rmfigure{figures/fig-29.png}{Fig.~29.}{Fig. 29: the lemniscate with chord $CS=z$ and its
semi-axis $CL$, from the original plate (Tav. III, Opere Matematiche vol. II).}

Secondo, che $\displaystyle\int \frac{-dz}{\sqrt{1-z^{4}}}$ esprime l'arco inverso $SL$ della
medesima curva, che corrisponde alla suddetta corda $CS$ ($z$).

Terzo, che $\displaystyle\int \frac{dz}{\sqrt{1+x^{4}}} = \frac{3}{2}\int dx\,\sqrt{1+x^{4}} -
\frac{1}{2}\sqrt{1+x^{4}}$, come si prova col calcolo.

Quarto, che $\displaystyle\int dx\,\sqrt{1+x^{4}}$ esprime l'arco $AQ$ (fig.~30) della parabola
cubica primaria corrispondente all'abscissa $AF = x$, purchè prendendo $a$ per l'unità, il
parametro di questa parabola sia $a\sqrt{3}$, e le ordinate sieno normali all'abscissa.

\rmfigure{figures/fig-30.png}{Fig.~30.}{Fig. 30: the primary cubic parabola with abscissa
$AF=x$, ordinate $FQ$, tangent portion $QV$, and origin $A$, from the original plate (Tav. III,
Opere Matematiche vol. II).}

Quinto, che $x\sqrt{1+x^{4}}$ denota la porzione $QV$ della tangente (fig.~30) compresa tra
l'ordinata $FQ$, e la retta indefinita $AV$, che dalla origine $A$ della parabola scorre parallela
alle ordinate.

Sesto che per conseguenza si avrà

\[ \int \frac{dx}{\sqrt{1+x^{4}}} = \frac{3}{2}\,\mathrm{arc.}\,AQ - \frac{1}{2}\,QV. \]

\emph{Esempj per mostrare il modo di applicare alla geometria i precedenti teoremi.}

\emph{Esempio I pel I teorema} (fig.~29, e 30). --- Nell'equazione (1) il segno arbitrario $\mp$
significhi il segno negativo, s'integri l'equazione (2) mediante il precedente scolio, e si
troverà

\[ \mathrm{Arc.}\,CS = \frac{3}{\sqrt{2}}\,\mathrm{arc.}\,AQ - \frac{1}{\sqrt{2}}\,QV. \]

Annullando $AF\,(x)$, si annullerà anche $CS\,(z)$, e però questa equazione è completa.

\emph{Esempio II pel II teorema} (fig.~29, e 30). --- Nell'equazione (3) il segno $\mp$
rappresenta il segno negativo, s'integri l'equazione (4), e si sco-
\origpage{307}
prirà

\[ \mathrm{Arc.}\,SL = \frac{3}{\sqrt{2}}\,\mathrm{arc.}\,AQ - \frac{1}{\sqrt{2}}\,QV. \]

La supposizione di $CS\,(z) = 1$, che annienta l'arco inverso $SL$, rende $AF\,(x) = 0$, e però
quest'equazione è completa.

\textbf{Scolio II.} --- Nel progresso di questo scritto io non proverò più, che l'equazioni sono
complete, dovendo bastare ai lettori il metodo, con cui mi sono regolato in questi due esempj per
accertarsi da loro medesimi della pienezza dell'equazioni seguenti; essi potranno ancora servirsi
dei quattro antecedenti teoremi per trarne altre maniere di misurare la lemniscata mediante
l'estensione della parabola cubica primaria, non ommettendo però, dov'è d'uopo, la sottrazione, o
l'aggiunta della quantità costante propria a rendere complete l'equazioni.

Potranno inoltre dedurre da questi teoremi delle verità affatto nuove, concernenti la
comparazione degli archi della suddetta parabola, ec.

Io mi contenterò di accennare, che siccome la misura della parabola cubica primaria dipende in
più maniere dall'estensione della lemniscata in virtù dei quattro antecedenti teoremi, e la
misura della lemniscata dipende dall'estensione dell'iperbola equilatera, e d'una specie di
elisse \emph{unitamente}, conforme ò scoperto nel I schediasma, ne segue, che la misura della
parabola cubica primaria dipende in più maniere dalla estensione delle due suddette sezioni
coniche \emph{unitamente}. Quest'invenzione non potrà non piacere a chi avrà considerato ciò, che
si legge negli Atti di Lipsia dell'anno 1695, pag. 64, e pag. 184 verso il fine.

\textbf{Teorema V.} --- Sieno le due equazioni infrascritte (9), e (10); io dico, che posta la
prima di esse sussiste anche l'altra

\[ \frac{u\sqrt{2}}{\sqrt{1-u^{4}}} = \frac{1}{z}\sqrt{1-\sqrt{1-z^{4}}}. \tag{9} \]

\[ \frac{dz}{\sqrt{1-z^{4}}} = \frac{2\,du}{\sqrt{1-u^{4}}}. \tag{10} \]

\textbf{Dimostrazione.} --- Ponendo nell'equazione (1) del I teorema in vece di $\mp$ il segno
negativo, e in cambio di $x$ il suo valore $\dfrac{u\sqrt{2}}{\sqrt{1-u^{4}}}$ notato
nell'equazione (5) del III teorema, si à l'equazione (9), indi sur-
\origpage{308}
rogando nell'equazione (2) del I teorema in luogo di $\dfrac{dx\sqrt{2}}{\sqrt{1+x^{4}}}$ il suo
valore $\dfrac{2\,du}{\sqrt{1-u^{4}}}$, che nasce dalla supposizione fatta di $x$ (come si vede,
moltiplicando per 2 l'equazione (6) del III teorema) ne viene l'equazione (10).

Il che dovea dimostrarsi.

\textbf{Corollario I} (fig.~29). --- Chiamisi $z$ la corda $CS$ della lemniscata, ed $u$ l'altra
corda $CI$, e integrando l'equazione (10), si troverà

\[ \mathrm{Arc.}\,CS = 2\,\mathrm{arc.}\,CI. \]

Laonde considerando la lettera $z$ come cognita nell'equazione (9), se ne trarrà il valore della
corda $CI\,(u)$ che taglia per mezzo l'arco diretto $CS$.

\textbf{Corollario II} (fig.~29). --- Dall'equazione (9) nasce quest'altra

\[ z = \frac{2u\sqrt{1-u^{4}}}{1+u^{4}}. \tag{11} \]

E questo valore di $z$ è tale, che chiamando $u$ la corda $CI$ considerata come cognita, e facendo
l'altra corda $CS$ eguale al secondo membro dell'equazione (11) l'arco diretto $CS$ è doppio
dell'altro arco diretto $CI$.

\textbf{Teorema IV.} --- Sieno le due equazioni infrascritte (12), e (13); io dico, che posta la
prima di esse, sussiste anche l'altra

\[ \frac{\sqrt{1-t^{4}}}{t\sqrt{2}} = \frac{1}{z}\sqrt{1-\sqrt{1-z^{4}}}. \tag{12} \]

\[ \frac{dz}{\sqrt{1-z^{4}}} = -\frac{2\,dt}{\sqrt{1-t^{4}}}. \tag{13} \]

\textbf{Dimostrazione.} --- Ponendo nell'equazione (1) del I teorema in vece di $\mp$ il segno
negativo, e in cambio di $x$ il suo valore $\dfrac{1}{t\sqrt{2}}\sqrt{1-t^{4}}$ notato
nell'equazione (7) del IV teorema, si à l'equazione (12). Sosti-
\origpage{309}
tuendo poi nell'equazione (2) del I teorema in luogo di $\dfrac{dx\sqrt{2}}{\sqrt{1+x^{4}}}$ il suo
valore $\dfrac{-2\,dt}{\sqrt{1-t^{4}}}$, che deriva dalla supposizione fatta di $x$ (come apparisce
moltiplicando per 2 l'equazione (8) del IV teorema) s'ottiene l'equazione (13).

Il che dovea dimostrarsi.

\textbf{Corollario I} (fig.~29). --- Si nomini $z$ la corda $CS$ della lemniscata, e ($t$)
l'altra corda $CH$, e ne risulterà arc. $CS = 2$ arc. $HL$.

\textbf{Corollario II} (fig.~29). --- Ordinando l'equazione (12), e ponendo in essa $u$ in cambio
di $z$, ed $r$ in vece di $t$, si giunge a quest'altra

\[ u = \frac{2r\sqrt{1-r^{4}}}{1+r^{4}}. \tag{14} \]

Di modo che nominando $r$ la corda $CO$, e assumendo l'altra corda $CI\,(u)$ eguale al secondo
membro dell'equazione (14) l'arco diretto $CI$ sarà doppio dell'arco inverso $OL$.

\textbf{Problema I.} --- Tagliare in tre parti uguali il quadrante della lemniscata.

\textbf{Soluzione} (fig.~29, e 31). --- In virtù del teorema VI, e suo I corollario l'arco
diretto $CS$ (fig.~29) sarà doppio dell'arco inverso $HL$, purchè chiamando $t$ la corda $CH$, e
$z$ l'altra corda $CS$ si abbia l'equazione (12). Suppongasi ora, che i due punti $S$, e $H$
coincidano in $T$, e questa supposizione renderà $CS\,(z) = CH\,(t)$, e l'arco diretto $CT$
(fig.~31) sarà doppio dell'arco inverso $TL$, che diverrà la terza parte del quadrante.

\rmfigure{figures/fig-31.png}{Fig.~31.}{Fig. 31: the lemniscate's quadrant with chord $CI$ and
points $T$, $L$ used in the trisection and quintisection problems, from the original plate
(Tav. III, Opere Matematiche vol. II).}

Ma supponendo $t = z$ nell'equazione (12), e poscia ordinandola nel debito modo, ec., si ottiene

\[ z = \sqrt[4]{-3+2\sqrt{3}}. \]

Laonde attribuendo alla corda $CT$ (fig.~31) questo valore di $z$, e pigliando l'arco diretto $CI$
eguale all'arco inverso $TL$ mediante il III teorema del I schediasma, si avrà lo scioglimento del
problema.

Il che dovea ritrovarsi.

\textbf{Problema II.} --- Tagliare in cinque parti uguali il quadrante della lemniscata.
\origpage{310}
\textbf{Soluzione} (fig.~29, e 31). --- Chiamisi $u$ la corda $CI$ (fig.~29) della lemniscata, e
$z$ l'altra corda $CS$, facciasi $z = \dfrac{2u\sqrt{1-u^{4}}}{1+u^{4}}$, e l'arco $CS$ sarà doppio
dell'arco $CI$ pel secondo corollario del V teorema. Si nomini poscia $r$ la corda $CO$ (fig.~29);
suppongasi la corda $CI\,(u) = \dfrac{2r\sqrt{1-r^{4}}}{1+r^{4}}$, e l'arco diretto $CI$ sarà
doppio dell'arco inverso $OL$ pel II corollario del VI teorema; dunque l'arco diretto $CS$ sarà
quadruplo dell'arco inverso $OL$.

Suppongasi ora $CS\,(z) = CO\,(r)$, i punti $S$, ed $O$ coincideranno in $T$, l'arco $TL$
(fig.~32) sarà un quinto del quadrante, l'arco $CI = 2$ arc. $TL$ conterrà due quinti dello stesso
quadrante, e l'arco $IT = $ arc. $CI$ ne conterrà gli altri due quinti.

Prendasi poscia mediante il III teorema del I schediasma l'arco inverso $EL$ (fig.~32) eguale
all'arco diretto $CI$, e l'arco diretto $CB$ eguale all'arco inverso $TL$, e i punti $B$, $I$,
$E$, $T$ taglieranno il quadrante in cinque parti eguali.

Resta solo a trovare il valore preciso della corda $CT$ (fig.~32) eguale nello stesso tempo a
$z$, e a $r$, e questo si troverà sostituendo nell'equazione (14) in luogo di $r$ la lettera $z$,
mentre si avrà

\[ u = \frac{2z\sqrt{1-z^{4}}}{1+z^{4}}. \tag{15} \]

E ponendo questo valore di $u$ nell'equazione (11) si giungerà a un'altra equazione, la quale non
conterrà altr'incognita che la lettera $z$, e dovrà reputarsi un'equazione dell'ottavo grado,
poichè se in essa si supporrà $z^{4} = b$, ne risulterà una nuova equazione appunto dell'ottavo
grado; il calcolo sarà più lungo, che difficile, purchè non si trascurino quelle maniere di
facilitarlo, che sono ben note ai periti analisti, e però stimo superfluo di stenderlo in questo
scritto. Avvertirò solamente, che nell'equazione espressiva del valore di $z^{4}$ la medesima
quantità di $z^{4}$ avrà due valori veri minori di $a^{4}$; il maggiore di questi due valori
esprimerà il quadratoquadrato di $CT\,(z)$, e il minore di essi esprimerà il quadratoquadrato di
$CI\,(u)$; imperciocchè se nell'equazione (15) si porrà invece di $z$ il suo valore in $u$ tratto
dall'equazione (11), si troverà una equazione affatto simile a quella, ch'esprime il valore di
$z^{4}$, ec.

Il che dovea ritrovarsi.

Il fu signor marchese de l'Hospital nel suo eccellente \emph{Trattato delle}
\origpage{311}
\emph{sezioni coniche}, lib. 9, prop. 9 insegna un modo di costruire l'equazioni dell'ottavo
grado.

\textbf{Teorema VII.} --- Sieno le due equazioni infrascritte (16), e (17); io dico, che posta la
prima di esse sussiste anche l'altra

\[ \frac{u\sqrt{2}}{\sqrt{1-u^{4}}} = \frac{\sqrt{1-z}}{\sqrt{1+z}}. \tag{16} \]

\[ \frac{-dz}{\sqrt{1-z^{4}}} = \frac{2\,du}{\sqrt{1-u^{4}}}. \tag{17} \]

Il che dovea dimostrarsi.

\textbf{Dimostrazione.} --- Nell'equazione (3) del II teorema il segno $\mp$ rappresenti il segno
negativo, e in vece di $x$ vi si ponga il suo valore $\dfrac{u\sqrt{2}}{\sqrt{1-u^{4}}}$ notato
nell'equazione (5) del III teorema, e si avrà l'equazione (16). Si surroghi poi nell'equazione (4)
del II teorema in luogo di $\dfrac{dx\sqrt{2}}{\sqrt{1+x^{4}}}$ il suo valore
$\dfrac{2\,du}{\sqrt{1-u^{4}}}$, che proviene dalla supposizione fatta di $x$ (come mostra
l'equazione (6) del III teorema moltiplicata per 2), e si avrà l'equazione (17).

Il che dovea dimostrarsi.

\textbf{Corollario} (fig.~29). --- Chiamando $z$ la corda $CS$ della lemniscata, ed $u$ l'altra
corda $CI$, e poscia integrando l'equazione (17), si ottiene Arc. $SL = 2$ arc. $CI$.

Purchè $u$, e $z$ abbiano il valore, che si deduce dall'equazione (16).

\textbf{Problema III} (fig.~33). --- Sia diviso per mezzo in $M$ il quadrante della lemniscata; e
sia dato l'arco diretto $CB$ minore della metà del quadrante: trovare l'arco intermedio $MI$
eguale all'arco diretto dato $CB$.

\rmfigure{figures/fig-33.png}{Fig.~33.}{Fig. 33: the lemniscate's quadrant with points $B$, $I$,
$M$, $E$ used in the bisection-transfer construction, from the original plate (Tav. III, Opere
Matematiche vol. II).}

\textbf{Soluzione.} --- I. Prendasi, mediante il teorema VII, e suo corollario, l'arco inverso
$EL$ doppio dell'arco diretto dato $CB$. II. Prendasi, mediante il teorema V, e suo I corollario,
l'arco diretto $CI$ suddoplo dell'arco diretto $CE$. Egli è visibile, che la somma degli archi
$CB + CI$
\origpage{312}
essendo eguale alla metà della somma degli archi $CE + EL$, sarà per conseguenza eguale alla metà
del quadrante; dunque

\[ \mathrm{arc.}\,CB + \mathrm{arc.}\,CI = \mathrm{arc.}\,CM, \]

e togliendo dall'uno, e l'altro membro di questa equazione l'arco comune $CI$, resta arc.
$CB = $ arc. $MI$.

Il che dovea dimostrarsi.

\textbf{Corollario} (fig.~33). --- Prendasi mediante il III teorema del I schediasma l'arco
inverso $FL$ eguale all'arco diretto $CI$, e si avrà

\[ \mathrm{arc.}\,Mf = \mathrm{arc.}\,MI = \mathrm{arc.}\,CB. \]

\textbf{Teorema VIII.} --- Sieno le due equazioni infrascritte (18), e (19); io dico, che posta la
prima di esse sussiste anche l'altra

\[ \frac{\sqrt{1-t^{4}}}{t\sqrt{2}} = \frac{\sqrt{1-z}}{\sqrt{1+z}}. \tag{18} \]

\[ \frac{-dz}{\sqrt{1-z^{4}}} = -\frac{2\,dt}{\sqrt{1-t^{4}}}. \tag{19} \]

\textbf{Dimostrazione.} --- Nell'equazione (3) del II teorema il segno $\mp$ rappresenti il segno
negativo, e in vece di $x$ vi si ponga il suo valore $\dfrac{\sqrt{1-t^{4}}}{t\sqrt{2}}$ espresso
nell'equazione (7) del IV teorema, e si avrà l'equazione (18). Sostituiscasi poscia nell'equazione
(4) del II teorema in cambio di $\dfrac{dx\sqrt{2}}{\sqrt{1+x^{4}}}$ il suo valore
$-\dfrac{2\,dt}{\sqrt{1-t^{4}}}$, che risulta dalla supposizione fatta di $x$ (conforme dimostra
l'equazione (8) del IV teorema moltiplicata per 2) e si giungerà all'equazione (19).

\textbf{Corollario} (fig.~29). --- Nominando $z$ la corda $CS$ della lemniscata, e ($t$) l'altra
corda $CH$, attribuendo a $z$, e a $t$ i loro debiti valori dedotti dall'equazione (18), e
integrando l'equazione (19) si scuopre l'arco inverso $SL$ eguale al doppio dell'arco inverso $HL$.

\textbf{Scolio III.} --- Dai teoremi V, VI, VII, e VIII, e loro corollarj potranno dedurre i
lettori dei modi di tagliare per mezzo il quadrante della lemniscata affatto uniformi alla maniera
di ciò fare da me scoperta
\origpage{313}
nel I schediasma, e questo servirà per prova della giustezza, e fecondità del mio metodo.

\textbf{Problema IV.} --- Posto, che il quadrante della lemniscata sia tagliato in un dato numero
di parti eguali, suddividere in due porzioni parimenti uguali ciascuna di dette parti.

\textbf{Soluzione.} --- Per maggior chiarezza, e brevità si chiamino archi diretti impari quegli
archi diretti della lemniscata, che contengono un numero impari delle parti eguali di essa. In
virtù del I corollario del V teorema si prendano gli archi diretti suddupli di tutti gli archi
impari; indi pel I corollario del III teorema del I schediasma si trovino gli archi inversi
eguali a tutti questi archi diretti suddupli, e si otterrà l'intento.

Il che dovea ritrovarsi.

\textbf{Scolio IV.} --- Dall'VIII teorema potranno dedurre i lettori un'altra maniera simile di
sciorre questo problema, che io per brevità tralascio.

\textbf{Corollario.} --- Poichè nel I schediasma ò insegnato il modo di tagliare il quadrante
della lemniscata in due parti eguali, e nel I, e II problema del presente scritto ò trovata la
maniera di segare il medesimo quadrante in tre, e cinque parti eguali; ne segue in vigore dello
scioglimento di questo problema, che il quadrante della lemniscata potrà dividersi
algebraicamente in tante parti eguali, quanti numeri si contengono in queste tre formole, cioè:
$2 \times 2^{m}$; $3 \times 2^{m}$; $5 \times 2^{m}$, nelle quali l'esponente $m$ significa
qualunque numero intiero positivo.

E questa è una nuova, e singolare proprietà della mia curva.

\end{document}
